\chapter{GamerServices}
\label{ch:gamerservices}

Real XNA's \cnans{GamerServices} namespace is Xbox-Live-only, and FNA --- targeting PC XNA
--- does not implement it at all. CNA goes further than FNA here on purpose: rather than
stub the whole namespace out, it built a real, locally-backed implementation held to the
actual Xbox~360 XNA~4.0 reference behavior as its fidelity bar, not to ``does nothing on a
PC.'' This chapter is organized around the resulting three-way split this namespace's own
documentation uses throughout: \textbf{Implemented} (real in-memory behavior matching the
reference), \textbf{Locally persisted} (real, disk-backed, but with no live online service
behind it), and \textbf{documented stub} (no local equivalent exists to build against at
all).

\section{Gamer, SignedInGamer, and a real aliasing hazard}

\cnaclass{Gamer} is the abstract base for \cnaclass{SignedInGamer}; \texttt{GetFromGamertag}
and \texttt{GetPartnerToken} (and their async pairs) are explicitly documented as ``not
supported in this platform's implementation'' --- the genuinely inert corners even this
locally-backed namespace did not attempt to reproduce, since they have no meaning without a
live Xbox Live service to query. One structural detail is worth knowing before writing any
code against this namespace: a \cnaclass{Gamer}'s \cnaclass{LeaderboardWriter}
(\S\ref{sec:leaderboards}) captures its owning \texttt{Gamer*} at construction, and because
neither \cnaclass{Gamer} nor \cnaclass{SignedInGamer} define custom copy or move
constructors, that captured pointer is copied verbatim by any ordinary copy or move of a
\cnaclass{Gamer} --- for instance, pushing a prvalue into a \texttt{std::vector}. Once a
\cnaclass{Gamer}'s \cnaclass{LeaderboardWriter} may actually be used, that \cnaclass{Gamer}'s
address must never move again; heap allocation is the documented, recommended way to avoid
this hazard entirely.

\cnaclass{SignedInGamer} is where the real, locally-persisted behavior lives:
\texttt{AwardAchievement} and \texttt{GetAchievements} are genuinely real, not stubs, backed
by a local persistence store. \texttt{IsFriend} always returns false, and
\texttt{GetFriends()} always returns an empty \cnaclass{FriendCollection} --- not because
\cnaclass{FriendGamer}/\cnaclass{FriendCollection} are unimplemented (they are fully real,
functioning classes), but because nothing anywhere in the codebase actually populates them;
this is the one honestly-acknowledged, still-open gap in this namespace's population story,
not a design limitation like the Xbox-Live-only methods above.

\section{Leaderboards: real persistence, CNA-original policy}
\label{sec:leaderboards}

\cnaclass{LeaderboardReader}/\texttt{Writer}/\texttt{Entry} are genuinely disk-backed --- a
deliberate expansion beyond FNA, whose own reference \cnaclass{LeaderboardReader}/%
\texttt{Writer} are identically all-\texttt{NotSupportedException} stubs, meaning there was
no real behavior to port here even if CNA had wanted to stay minimal. Because no reference
sort or paging policy existed to match, CNA's own policy is explicitly original rather than
derived: entries sort by rating descending; paging around a specific pivot gamer falls back
to the top of the board if that gamer has no entry at all; and a persisted gamertag with no
currently signed-in matching \cnaclass{Gamer*} is simply skipped, since a
\cnaclass{LeaderboardEntry} needs a live gamer object to attach itself to. Every
\texttt{Begin}/\texttt{End} async pair on this class completes synchronously, matching the
project-wide pattern already seen for \cnaclass{NetworkSession} in
Chapter~\ref{ch:networking}. One genuinely subtle, load-bearing detail: real XNA's own
\cnaclass{LeaderboardWriter} has no explicit submit or commit method at all, so CNA treats
assigning \texttt{Rating} as the closest honest equivalent of ``I'm finished, persist this''
--- but only for an entry actually obtained through a \cnaclass{LeaderboardWriter}, since only
those have a persistence hook installed; setting \texttt{Rating} on an entry obtained from a
\cnaclass{LeaderboardReader} is deliberately a no-op.

\section{Achievements}

\cnaclass{Achievement} and \cnaclass{AchievementCollection} are real and locally persisted,
following the same pattern as leaderboards. One method is a deliberate, faithful exception to
that: \texttt{Achievement::GetPicture()} always throws \texttt{NotImplementedException},
matching FNA's own unimplemented stub --- real Xbox~360 achievement artwork streamed from
Xbox Live on request and has no meaningful local substitute, so this is treated the same way
this book's other chapters treat a genuinely unreachable capability: named and left throwing,
not quietly worked around.

\section{Guide: mostly a stub, with two real, CNA-built exceptions}

\cnaclass{Guide} is where this namespace's ``real versus stub'' distinction is sharpest. Most
of its surface --- \texttt{ShowFriends}, \texttt{ShowMarketplace}, \texttt{ShowGamerCard},
\texttt{ShowSignIn}, \texttt{ShowParty}, and the rest of the \texttt{Show*} family --- is a
genuine, documented no-op: there is no system UI to show for any of them, matching FNA's own
minimal PC-side stub. Two specific methods are the deliberate exception, and are real, CNA-built
overlay UI, not stubs at all: \texttt{BeginShowMessageBox}/\texttt{EndShowMessageBox} and
\texttt{BeginShowKeyboardInput}/\texttt{EndShowKeyboardInput}. Unlike almost everything else
in this namespace, these do \emph{not} complete synchronously --- they remain genuinely
pending until the game's own \texttt{Draw()} loop calls a dedicated render hook once per
frame, which draws the message box or keyboard prompt using a caller-supplied
\cnaclass{SpriteBatch}/\cnaclass{SpriteFont}, and polls real keyboard/mouse input to detect a
click, Enter, or Escape. The keyboard-input path captures real UTF-16 text safely across
surrogate pairs, and its password-mode masking only affects the \emph{displayed} text --- the
value actually returned is always the real typed text, never the masked version.
\texttt{Guide.IsVisible} is correspondingly real (not the fixed ``always false'' constant FNA
itself uses off Xbox) --- true exactly when one of these two real overlays is currently
pending.

\section{Dispatcher and component}

\cnaclass{GamerServicesComponent} must be added to \texttt{Game.Components}
(this book's game-loop chapter) before any \cnaclass{Gamer}/\cnaclass{Guide} API is
used, matching real XNA~4.0's own PC requirement exactly. \cnaclass{GamerServicesDispatcher}
is the actual entry point creating and freeing the local, stub-style
\cnaclass{SignedInGamer} set each time it initializes, and pumping the async machinery every
other class in this namespace relies on for its synchronously-completing
\texttt{Begin}/\texttt{End} pairs.

\section{Summary table}

\begin{center}
\begin{tabular}{ll}
\toprule
\textbf{Feature} & \textbf{Status} \\
\midrule
\cnaclass{SignedInGamer}/\cnaclass{GamerCollection} & Implemented \\
Achievements & Locally persisted \\
Leaderboards & Locally persisted, CNA-original policy \\
\cnaclass{GamerPresence}/\cnaclass{GamerPrivileges} & Implemented \\
\cnaclass{FriendCollection}/\cnaclass{FriendGamer} & Real classes, population always empty \\
\texttt{Guide.Show*} (most) & Documented no-op stub \\
\texttt{Guide.BeginShowMessageBox}/\texttt{KeyboardInput} & Implemented, CNA-built overlay UI \\
Avatar (base API) & Faithfully inert --- see Chapter~\ref{ch:avatar} \\
\bottomrule
\end{tabular}
\end{center}
